% ----------------------------------------------------------------------
\subsection{FC3 Compression}
\label{sub:fc3_compression}
% ----------------------------------------------------------------------

Storing the full FC3 tensor requires $\cO(\nprim\,\ndof^{2}) = \cO(\Ncell^{2}\nprim^{3})$
memory, which is prohibitive for large-scale simulations. Beyond storage, the dense four-tensor
contraction entering the phonon--phonon self-energy scales as $\cO(\ndof^{4})$ per frequency
point and dominates the cost of the quantum-transport simulation. It is hence desirable to
efficiently approximate the FC3 interactions.

Effectively every published anharmonic phonon NEGF implementation avoids the problem by
dropping self-energy and FC3 entries at large atomic separations
\cite{guo_quantum_2020, guo_anharmonic_2021, luisier_atomistic_2012, li_anharmonic_2009,
      polanco_nonequilibrium_2021, mingo_anharmonic_2006}.
The standard choice is a real-space pair-distance cut-off: keep only those triplets
$(0\kappa, l'\kappa', l''\kappa'')$ whose largest pairwise distance lies below some
$r_\mathrm{c}$. The same idea is used in the BTE workflows of ShengBTE and its successors,
which restrict FC3 to a few nearest-neighbour shells \cite{li_shengbte_2014}.

The drawback is that the approximation error is uncontrolled: there is no variational bound
on what is discarded, and the truncation removes genuine non-local contributions that can
matter quantitatively in the phonon self-energy. In anharmonic NEGF this is already visible in
silicon. \citeauthor{guo_quantum_2020} use first-principles FC3 in a quantum-transport treatment
of cross-plane heat flow in Si thin films, where they compare first- and third-nearest-neighbour
FC3 cutoffs \cite{guo_quantum_2020}. Using the same FC3 data in ShengBTE, they report bulk Si
thermal conductivities at 300~K of $120.69$, $136.46$, and $147.13~\mathrm{W\,m^{-1}K^{-1}}$ when
FC3 was truncated at the first, second, and third neighbour shell, respectively, so that a
first-shell cut-off underestimates the converged bulk value by nearly $20\%$. For bulk silicon,
including up to 4th nearest neighbour interactions gives converged results
\cite{li_shengbte_2014, qin_accelerating_2018}.

For three-dimensional covalent crystals, the cut-off problem is usually less severe than in
two-dimensional materials. \citeauthor{mcgaughey_phonon_2019} report heat conductivities
$\kappa_\ell(300\,\mathrm{K}) = 150, 152, 155, 151,$ and $149~\mathrm{W\,m^{-1}K^{-1}}$ for
cubic cut-offs of $2$, $3$, $4$, $5$, and $6$ neighbour shells, respectively in bulk Si, indicating
that once a few shells are retained the residual cut-off dependence is already small
\cite{mcgaughey_phonon_2019}. \citeauthor{jiang_how_2022} reached the same qualitative conclusion
by contrasting graphene and silicene against bulk Si: whereas the 2D systems require inclusion of
up to 14th order NN shells, bulk Si shows ``no convergence issue'' even for very small neighbour
cut-offs \cite{jiang_how_2022}.

That behaviour is, however, not universal across 3D materials. \citeauthor{qin_accelerating_2018}
showed that bulk SnSe still exhibits strong interactions at distances around $6.15~\text{\AA}$,
and that these coincide with a sharp change in the converged $\kappa_\ell$ once the pair cut-off
exceeds roughly $6~\text{\AA}$ \cite{qin_accelerating_2018}. They also show that for the 2D
materials studied, the observables in BTE diverge for increasing $q$-grid when the cut-off is too
small.

We therefore seek a representation of the FC3 tensor that
\begin{itemize}
    \item provides a controllable approximation error,
    \item achieves significant compression of the stored data, and
    \item turns the dense $\cO(\ndof^{4})$ self-energy contraction into a sum of low-order
          contractions whose cost scales polynomially in a parameter~$R$.
\end{itemize}
A large body of work has developed such representations in similar settings.
In quantum chemistry, tensor hypercontraction (THC) \cite{hohenstein_tensor_2012,
parrish_tensor_2012} and density fitting \cite{whitten_coulombic_1973} factorise the four-index
electron-repulsion integral tensor into products of matrices. In many-body physics, tensor
networks (matrix product states, tensor trains) \cite{oseledets_tensor-train_2011,
cichocki_tensor_2016} compress high-dimensional objects at controlled error. For the present
setting, lattice dynamics, tensor decomposition has recently proved effective:
\citeauthor{luo_data-driven_2024} used a truncated SVD on the electron-phonon matrix in Wannier
basis and recovered high-accuracy transport, spin, and superconducting observables from $1$--$2\%$
of the singular values \cite{luo_data-driven_2024}. The same group then proposed the permanent
CP (PCP) ansatz for the $n$-IFC tensors and reported $10^{3}$--$10^{4}$ compression factors
across Si, HgTe, MgO, TiNiSn and monoclinic ZrO$_2$ while retaining $>98\%$ accuracy on
$\kappa_\ell$ for three- and four-phonon scattering \cite{luo_tensor_2025}. Compressive-sensing
lattice dynamics (CSLD) is a method applied to a related problem at the fitting stage: an
$\ell_1$-regularised regression on DFT forces selects a sparse set of non-negligible FC entries
directly \cite{zhou_lattice_2014, zhou_compressive_2019}. Projector-based approaches, implemented
in the symfc software, construct a complete orthonormal basis for FCs compatible with all
crystal symmetries \cite{seko_projector-based_2024}, while group-theoretical reductions enumerate
the independent FC3 entries from the space group \cite{fu_group_2019}. In the following, we take
the fully-populated FC3 tensor as given (from finite displacements, DFPT, or CSLD) and compress
its representation. However, one can likely also fit a compressed tensor from DFT data with an
adaption of CSLD or similar methods without ever obtaining a full FC3 tensor.

\citet{kolda_tensor_2009} organise tensor compression into two principal families: the
Tucker decomposition, which writes a tensor as a core tensor multiplied by a factor matrix
along each mode (generalising principal-component analysis), and the CP
(CANDECOMP/PARAFAC) decomposition, which writes it as a sum of rank-one outer products
(generalising the matrix SVD in a different direction). Applied to a tensor with the FC3 index
pattern and combined with the possible symmetry constraints on the factor matrices, these two
families yield five canonical ans\"atze, ordered here from least to most
tensor-structure-aware:
\begin{enumerate}
  \item Truncated matricization SVD, i.e.\ the classical matrix SVD applied to an matricization of
        the tensor. fully algebraic and Eckart--Young optimal in the chosen matricization, but
        separating only one leg from the other two.
  \item Truncated HOSVD (Tucker), with independent factor matrices per mode and a small dense core.
  \item CP, obtained by restricting the Tucker core to be diagonal, with three
        independent rank-$R$ factor matrices and no symmetry constraint.
  \item INDSCAL \cite{carroll_analysis_1970}, obtained from CP by setting the
        factor matrices on the two internal legs equal, exploiting the
        $S_2$ symmetry $\Phi_{\mu j k} = \Phi_{\mu k j}$.
  \item Symmetric CP, equivalent to the Waring decomposition of a homogeneous
        cubic polynomial \cite{brachat_symmetric_2010, comon_symmetric_2008,
        oeding_eigenvectors_2013}, obtained from CP by setting all three factor
        matrices equal (using the $p(\mu)$ lift of a primitive-cell DOF to its
        supercell counterpart in cell zero), exploiting the full $S_3$
        symmetry.
\end{enumerate}
Entries 1 and 2 are non-iterative. Entries 3--5 are NP-hard in the worst case
\cite{hillar_most_2013} and use iterative optimisation, with the known caveats of
rank-degeneracy and non-closedness of rank-$R$ sets
\cite{stegeman_degeneracy_2006, de_silva_tensor_2008}. Each added symmetry constraint from
entry~2 onwards reduces the parameter count at matched rank while shrinking the set of
representable tensors. As discussed in %\cref{sub:fc3_observable_error}
\todo{write section about relation between error and observable}
, the most-constrained ansatz is not necessarily the least accurate per
parameter.

% ----------------------------------------------------------------------
\subsection{Truncated matricization SVD}
\label{sub:fc3_matsvd}
% ----------------------------------------------------------------------
Stack the $3\nprim$ slices $M_\mu = \Phi_{\mu,:,:}\in\mathbb{R}^{\ndof\times\ndof}$
into a single matrix
\begin{equation}\label{eq:fc3_matsvd_stack}
  M_{\mathrm{stack}}
  = \begin{pmatrix} M_1 \\ \vdots \\ M_{3\nprim} \end{pmatrix}
  \in \mathbb{R}^{3\nprim\,\ndof \,\times\, \ndof}.
\end{equation}
A truncated matrix SVD at rank $R$ gives
\begin{equation}\label{eq:fc3_matsvd_ansatz}
    \tilde\Phi_{\mu j k}^{\mathrm{mSVD}}
    = \sum_{r=1}^{R} F_r(\mu,j)\; h_r(k),
\end{equation}
with $F_r(\mu,j) = \sigma_r\,[\vc{u}_r]_{(\mu-1)\ndof+j}$ and $h_r(k)=[\vc{v}_r]_k$, where
$\sigma_r$, $\vc{u}_r$, $\vc{v}_r$ are the $R$ leading singular values, left singular vectors,
and right singular vectors of $M_{\mathrm{stack}}$. This is equivalent to a Tucker1 decomposition
with compression on mode~3 only \cite[\S4]{kolda_tensor_2009}.

\paragraph{Computation.} A single matrix SVD of $M_{\mathrm{stack}}$, computed by
standard dense LAPACK routines or, for large $\ndof$, by randomised linear algebra algorithms
\cite{halko_finding_2010}. The Eckart--Young--Mirsky theorem
\cite{eckart_approximation_1936, mirsky_symmetric_1960} guarantees that, at the chosen $(\mu,j)|k$ bipartition,
the truncation is globally optimal among all rank-$R$ approximations in the Frobenius norm.
This makes it the natural baseline against which the structurally richer decompositions of
\cref{sub:fc3_tucker,sub:fc3_cp,sub:fc3_indscal,sub:fc3_symcp} must be compared. The ansatz does
not exploit the $S_2$ symmetry on $(j,k)$: modes~1 and~2 are folded into a single flat index, so
parameters are wasted on storing redundant information and the reconstruction is not exactly
symmetric in $(j,k)$ unless restored in post-processing.

\paragraph{Parameter count.}
\begin{equation}\label{eq:fc3_matsvd_params}
  N_{\mathrm{param}}^{\mathrm{mSVD}}(R)
  = R\bigl((3\nprim+1)\,\ndof + 1\bigr)
  \;\sim\; \cO\!\bigl(R\,\nprim\,\ndof\bigr)
  \;=\; \cO\!\bigl(R\,\Ncell\,\nprim^{2}\bigr),
\end{equation}
one factor of $\nprim$ worse than the CP-family ans\"atze below.

\paragraph{Self-energy kernel.} The separable structure factorises the ring contraction entering
the phonon--phonon self-energy into two independent bilinear forms. The $\vq$-convolution is an
FFT over the $R^{2}$ rank pairs, replacing the dense $\ndof^{4}$ contraction by
$R^{2}\ndof$-scaling accumulations.

% ----------------------------------------------------------------------
\subsection{Truncated HOSVD (Tucker decomposition)}
\label{sub:fc3_tucker}
% ----------------------------------------------------------------------
The Tucker ansatz writes FC3 as a multilinear contraction of a small core tensor with one factor
matrix per mode \cite{tucker_mathematical_1966, kolda_tensor_2009}:
\begin{equation}\label{eq:fc3_tucker_ansatz}
    \tilde\Phi_{\mu j k}^{\mathrm{HOSVD}}
    = \sum_{p=1}^{R_1}\sum_{q=1}^{R_2}\sum_{r=1}^{R_2}
      G_{p q r}\; A_{\mu p}\; B_{j q}\; B_{k r},
\end{equation}
with $A\in\mathbb{R}^{3\nprim\times R_1}$, $B\in\mathbb{R}^{\ndof\times R_2}$ columnwise
orthonormal, and a core $G\in\mathbb{R}^{R_1\times R_2\times R_2}$. To preserve the $S_2$
symmetry of the two internal legs, the factor matrices on modes~2 and~3 can be chosen identical ($B=C$)
and the core is constrained to satisfy $G_{pqr}=G_{prq}$.

\paragraph{Computation.} The truncated higher-order SVD (HOSVD) of
\citet{de_lathauwer_multilinear_2000} produces $A$ and $B$ in closed form as the $R_1$ and $R_2$
leading left singular vectors of the mode-$1$ and mode-$2$ unfoldings of $\Phi$, followed by the
core contraction $G = \Phi\times_1 A^\top\times_2 B^\top\times_3 B^\top$ (which is automatically
$S_2$-symmetric when $\Phi$ is). This is quasi-optimal: the truncation error is bounded by a
factor $\sqrt{3}$ above the best rank-$(R_1,R_2,R_2)$ approximation error
\cite{hackbusch_tensor_2019, oseledets_tensor-train_2011}. The fit can optionally be refined by a
few steps of higher-order orthogonal iteration (HOOI), an alternating least squares (ALS)-type loop
 \cite{de_lathauwer_best_2000}. The special case $R_1=3\nprim$, $A=I$ (compression of the two
internal legs only) is the Tucker1 baseline and is obtained by a single SVD of the corresponding
unfolding.

\paragraph{Parameter count.}
\begin{equation}\label{eq:fc3_tucker_params}
  N_{\mathrm{param}}^{\mathrm{HOSVD}}(R_1,R_2)
  = 3\nprim\,R_1 + \ndof\,R_2 + R_1\,\tfrac{R_2(R_2+1)}{2}
  \;\sim\; \cO\!\bigl(R_2\,\ndof + R_1\,R_2^{2}\bigr),
\end{equation}
where the $S_2$ constraint halves the nominal core cost $R_1 R_2^{2}$ to $R_1 R_2(R_2+1)/2$.
The $R_2^{2}$-scaling core term makes Tucker the most parameter-heavy of the five ans\"atze
when $R_2$ is not small compared to $\ndof^{1/2}$. for strongly anharmonic systems where the
effective rank is moderate this is mitigated by allowing $R_1\ll R_2$.

\paragraph{Self-energy kernel.} The self-energy contraction factorises into three mode-wise
contractions against $A$ and $B$, leaving the $R_1\times R_2\times R_2$ core as a small dense
sub-problem. The total cost scales as $R_2^{2}(R_1+\ndof)$ per frequency point rather than
$\ndof^{4}$.

% ----------------------------------------------------------------------
\subsection{Canonical Polyadic decomposition (CP)}
\label{sub:fc3_cp}
% ----------------------------------------------------------------------
Dropping the off-diagonal entries of the Tucker core yields the CP approximation, also
known as CANDECOMP and PARAFAC \cite{harshman_foundations_1970, hitchcock_expression_1927, carroll_analysis_1970}, a sum of
rank-one outer products with three independent factor vectors per term:
\begin{equation}\label{eq:fc3_cp_ansatz}
    \tilde\Phi_{\mu j k}^{\mathrm{CP}}
    = \sum_{r=1}^{R} \lambda_r\;a_r(\mu)\;b_r(j)\;c_r(k).
\end{equation}
No symmetry between the modes is imposed. For a tensor that is
$S_2$- or $S_3$-symmetric, the fit does not necessarily converge to a symmetric
solution \cite{ten_berge_typical_2004}. Uniqueness can be shown under certain conditions
\cite{kruskal_three-way_1977, sidiropoulos_uniqueness_2000} and is generically
guaranteed for the FC3 tensor of interest here
\cite{chiantini_generic_2012}.

\paragraph{Computation.} The standard algorithm is alternating least squares
(CP-ALS), which cycles through the three modes solving, at each step, a
linear least-squares problem for one factor matrix while the other two are
fixed \cite[Sec.~3.4]{kolda_tensor_2009}. Convergence is only to a
stationary point and can be slow. Variants include
damped Gauss--Newton (dGN), PMF3 \cite{paatero_weighted_1997},
line-search extrapolation (ELS) \cite{rajih_enhanced_2008}, and
simultaneous-diagonalisation approaches \cite{de_lathauwer_link_2006}
that recast CP as a generalised eigenproblem when
$R\le\min(3\nprim,\ndof)$.

\paragraph{Parameter count.}
\begin{equation}\label{eq:fc3_cp_params}
  N_{\mathrm{param}}^{\mathrm{CP}}(R)
  = R\,\bigl(3\nprim + 2\,\ndof + 1\bigr)
  \;\sim\; \cO\!\bigl(R\,\ndof\bigr)
  \;=\; \cO\!\bigl(R\,\Ncell\,\nprim\bigr).
\end{equation}

% ----------------------------------------------------------------------
\subsection{INDSCAL (CP with internal-leg symmetry)}
\label{sub:fc3_indscal}
% ----------------------------------------------------------------------
Individual-differences-in-scaling (INDSCAL) \cite{carroll_analysis_1970, kolda_tensor_2009} is the
specialisation of CP in which the factor matrices on two modes are constrained equal, tailored
to tensors whose slices $\Phi_{\mu,:,:}$ are symmetric:
\begin{equation}\label{eq:fc3_indscal_ansatz}
    \tilde\Phi_{\mu j k}^{\mathrm{INDSCAL}}
    = \sum_{r=1}^{R} d_r(\mu)\; v_r(j)\; v_r(k),
\end{equation}
with $d_r\in\mathbb{R}^{3\nprim}$ and $v_r\in\mathbb{R}^{\ndof}$. This imposes the $S_2$
symmetry $\Phi_{\mu j k}=\Phi_{\mu k j}$. Typical and maximal ranks for
partially symmetric tensors were characterised by
\citet{ten_berge_typical_2004}: for the shapes relevant here, they coincide with or
are strictly smaller than the ranks of generic nonsymmetric tensors of the same shape, so the
symmetry constraint is not costly in expressive power.

\paragraph{Computation.} \citet{carroll_analysis_1970} originally proposed running unconstrained
CP-ALS with the two internal factors $B$ and $C$ updated separately, relying on symmetry of the
data to drive $B\to C$. \citet{ten_berge_typical_2004, ten_berge_carroll_2009} showed that this convergence is generic
but not guaranteed. Modern implementations impose equality explicitly via a symmetric ALS or via
a Gauss--Newton method \cite{kolda_tensor_2009, kolda_numerical_2015}.
A cheap closed-form alternative performs a matrix SVD of $\Phi$ reshaped to
$3\nprim\times\ndof^{2}$ followed by an eigendecomposition of each symmetric slice of the
resulting right factors, and keeps the $R$ terms with largest amplitude. This is not jointly
optimal (the total rank needed to match a given error is typically larger than that of a
jointly-fitted INDSCAL) but is useful as an initialisation.

\paragraph{Parameter count.}
\begin{equation}\label{eq:fc3_indscal_params}
  N_{\mathrm{param}}^{\mathrm{INDSCAL}}(R)
  = R\,\bigl(3\nprim + \ndof\bigr)
  \;\sim\; \cO\!\bigl(R\,\ndof\bigr)
  \;=\; \cO\!\bigl(R\,\Ncell\,\nprim\bigr),
\end{equation}
one factor of $\ndof$ smaller than CP at matched rank, from merging $b_r\equiv c_r$.

\paragraph{Self-energy kernel.} Identical in scaling to CP, $\cO(R^{2}\ndof)$ per frequency, but
with the additional property that the reconstruction is exactly $S_2$-symmetric in $(j,k)$, so
ASR-sensitive observables receive only the quadratic-in-error contribution discussed in
\todo{actually discuss this}%\cref{sub:fc3_observable_error}.


% ----------------------------------------------------------------------
\subsection{Symmetric CP (Waring) decomposition}
\label{sub:fc3_symcp}
% ----------------------------------------------------------------------
The symmetric CP decomposition imposes equality of all three factor matrices, $A=B=C$, on the
fully $S_3$-symmetrised FC3 tensor obtained via the supercell lift
$p:\{1,\dots,3\nprim\}\to\{1,\dots,\ndof\}$. The Waring parametrisation reads:
\begin{equation}\label{eq:fc3_waring_ansatz}
    \tilde\Phi_{\mu j k}^{\mathrm{Waring}}
    = \sum_{r=1}^{R} \lambda_r\;
      v_r\!\bigl(p(\mu)\bigr)\; v_r(j)\; v_r(k),
\end{equation}
with $v_r\in\mathbb{R}^{\ndof}$ and $\lambda_r\in\mathbb{R}$. This is the tensor form of the
classical Waring problem: expressing a homogeneous cubic polynomial as a sum of $R$ cubes of
linear forms \cite{brachat_symmetric_2010, comon_symmetric_2008, oeding_eigenvectors_2013}. For
generic third-order symmetric tensors on $\ndof$ indices, the Waring rank equals
$\lceil\binom{\ndof+2}{3}/\ndof\rceil$ by the Alexander--Hirschowitz theorem
\cite{alexander_generic_1997} (which is a non-trivial statement and proven by Alexander and Hirschowitz
in a set of 5 papers of over 100 pages). An alternative ``permanent-CP'' parametrisation, introduced
by \citet{luo_tensor_2025} as PCP, uses three distinct vectors per term explicitly symmetrised
over $S_3$,
\begin{equation}\label{eq:fc3_pcp_ansatz}
    \tilde\Phi_{\mu j k}^{\mathrm{PCP}}
    = \sum_{\xi=1}^{R} \frac{\lambda_\xi}{6}
      \sum_{\sigma\in S_3}
      u^{\xi}_{\sigma_1}\!\bigl(p(\mu)\bigr)\;
      u^{\xi}_{\sigma_2}(j)\;
      u^{\xi}_{\sigma_3}(k),
\end{equation}
with $u^{\xi}_1, u^{\xi}_2, u^{\xi}_3\in\mathbb{R}^{\ndof}$. PCP fits the same class of symmetric
tensors as Waring but with more capacity per term (at threefold parameter cost) and a smoother
optimisation landscape.

\paragraph{Computation.} Algorithms fall into two categories. Algebraic methods reconstruct the
decomposition from moment and flattening matrices via generalised eigenproblems and are exact up
to numerical precision when the rank is below the generic bound
\cite{brachat_symmetric_2010, oeding_eigenvectors_2013}. Numerical methods optimise
$\|\Phi-\tilde\Phi\|_F^{2}$ directly: the shifted symmetric higher-order power method (SS-HOPM)
\cite{kolda_shifted_2011} and the subspace power method \cite{kileel_subspace_2025} extract
components one at a time with provable convergence for low-rank tensors; all-at-once gradient
and Gauss--Newton methods \cite{kolda_numerical_2015} fit all $R$ components simultaneously.
Degeneracy and diverging-norm instabilities can appear because the set of symmetric tensors of
rank $\le R$ is not closed in general
\cite{stegeman_degeneracy_2006, de_silva_tensor_2008}. Remedies include norm penalties and
second-order line searches.

\paragraph{Parameter count.}
\begin{align}
  N_{\mathrm{param}}^{\mathrm{Waring}}(R) &= R\,(\ndof + 1)
    \;\sim\; \cO\!\bigl(R\,\ndof\bigr)
    \;=\; \cO\!\bigl(R\,\Ncell\,\nprim\bigr),
    \label{eq:fc3_waring_params}\\
  N_{\mathrm{param}}^{\mathrm{PCP}}(R) &= R\,(3\,\ndof + 1)
    \;\sim\; \cO\!\bigl(R\,\ndof\bigr)
    \;=\; \cO\!\bigl(R\,\Ncell\,\nprim\bigr).
    \label{eq:fc3_pcp_params}
\end{align}

\paragraph{Self-energy kernel.} Identical in scaling to CP and INDSCAL, $\cO(R^{2}\ndof)$ per
frequency, but with the additional property that the reconstruction is exactly $S_3$-symmetric
in $(p(\mu),j,k)$, so ASR-sensitive observables receive only the quadratic-in-error contribution
discussed in \todo{discuss}
%\cref{sub:fc3_observable_error}.


% ----------------------------------------------------------------------
\subsection{Further decompositions and structured variants}
\label{sub:fc3_other}
% ----------------------------------------------------------------------
The five ans\"atze of
\cref{sub:fc3_matsvd,sub:fc3_tucker,sub:fc3_cp,sub:fc3_indscal,sub:fc3_symcp} do not exhaust the
useful choices. Further approximations appear in the broader literature and can either be
combined with them or replace them in specific regimes.

\paragraph{Cross / CUR / skeleton approximations.} Rather than fitting a compressed form to a
precomputed $\Phi$, one can build the approximation from a small number of sampled fibers of the
tensor and interpolate the rest:
\begin{equation}
  \tilde\Phi_{\mu j k}^{\mathrm{CUR}}
  = \sum_{p,q,r}\Phi_{\mu,J_q,K_r}\,
    \Phi_{I_p,J_q,K_r}^{-1}\,\Phi_{I_p,j,k},
\end{equation}
where $I,J,K$ are index subsets of size $\sim R$ selected to maximise the volume of the
intersection subtensor, and $\Phi_{I,J,K}$ is inverted in the appropriate pseudoinverse sense.
This bypasses the construction of the full FC3 altogether, which is normally the computational
bottleneck (finite-displacement DFT). Rank-$R$ tensor cross
\cite{oseledets_tt-cross_2010, savostyanov_quasioptimality_2014} accesses $\cO(R^{2}\ndof)$ entries.
Adaptive cross approximation \cite{bebendorf_approximation_2000} and CUR-style subset selection
\cite{goreinov_theory_1997, mahoney_cur_2009} give analogous matrix-level
guarantees, and fiber-sampling Tucker (FSTD) of \citet{caiafa_generalizing_2010} yields an explicit
Tucker form from $\cO(R^{N-1})$ fibers. The tensor-CURT extension of
\citet{song_relative_2019} provides the first relative-error bounds in this setting. For
FC3 the main challenge is that cross methods have no intrinsic notion of ASR or $S_3$ symmetry.
A hybrid that samples in the symmetric CP parameter space would combine the two sets of
guarantees.

\paragraph{Tensor train and hierarchical Tucker.} The tensor-train (TT) format
\cite{oseledets_tensor-train_2011} writes an $N$-index tensor as a contraction of $N$ three-index
cores with bond ranks $(R_1,\dots,R_{N-1})$. For third-order tensors this collapses to
\begin{equation}
  \tilde\Phi_{\mu j k}^{\mathrm{TT}}
  = \sum_{\alpha_1,\alpha_2} G_1(\mu,\alpha_1)\,
    G_2(\alpha_1,j,\alpha_2)\,G_3(\alpha_2,k),
\end{equation}
which is structurally a Tucker1 SVD of the $(\mu)|(j,k)$ unfolding followed by a Tucker1 SVD of
the $(\alpha_1 j)|(k)$ matricisation of the residual. It does not offer independent advantages
beyond \cref{sub:fc3_tucker}. TT and the hierarchical Tucker (HT) format of
\citet{hackbusch_tensor_2019} become first-class methods at order
$\ge 4$ and are the natural representation for FC4 (four-phonon) tensors, where CP and Tucker
lose efficiency. \citet{lee_solving_2026} recently used Matrix Product States (= TT) to solve the
Peierls--Boltzmann equation in silicon.

\paragraph{Block-term decomposition (BTD).} \citet{de_lathauwer_decompositions_2008, de_lathauwer_decompositions_2008-1}
generalise both CP and Tucker by writing the tensor as a sum of low-multilinear-rank Tucker blocks,
\begin{equation}
  \tilde\Phi_{\mu j k}^{\mathrm{BTD}}
  = \sum_{r=1}^{R}\bigl(\mathcal{G}_r\times_1 A_r\times_2 B_r\times_3 C_r\bigr)_{\mu j k},
\end{equation}
with each $\mathcal{G}_r$ of small core shape (e.g.\ the rank-$(L,L,1)$ BTD has
$\mathcal{G}_r$ of shape $L\times L\times 1$). CP is the rank-$(1,1,1)$ special case; Tucker is
the single-block special case. For FC3 this structurally expresses block patterns like
per-atom-species or per-shell contributions in a single decomposition, and inherits the
uniqueness and ALS algorithms of the parent families.

\paragraph{Constrained / symmetry-adapted decompositions (CANDELINC).} CANDELINC
\cite{douglas_carroll_candelinc_1980} is CP with factor matrices constrained to user-chosen
linear subspaces,
\begin{equation}
  A=\Phi_A\hat A,\quad B=\Phi_B\hat B,\quad C=\Phi_C\hat C,
\end{equation}
where $\Phi_A,\Phi_B,\Phi_C$ are columnwise-orthonormal basis matrices. It reduces to an ordinary
CP fit in the reduced coordinates $(\hat A,\hat B,\hat C)$ applied to the projected tensor
$\hat\Phi = \Phi\times_1\Phi_A^\top\times_2\Phi_B^\top\times_3\Phi_C^\top$. For FC3 this is the
natural framework to couple any of the five ans\"atze with a space-group-compatible basis:
choose $\Phi_A,\Phi_B,\Phi_C$ to span the ASR-respecting, symmetry-invariant subspace (symfc
\cite{seko_projector-based_2024}, Fu--Marianetti \cite{fu_group_2019}, CSLD
\cite{zhou_lattice_2014, zhou_compressive_2019}), then fit INDSCAL or symmetric CP on the reduced
tensor.

\paragraph{Sparse / structured CP.} Imposing sparsity or smoothness priors on the CP factor
vectors yields decompositions whose rank-$R$ components are localised on a few atoms or DOFs.
As an example, one can penalise the loss with
\begin{equation}
  L(A,B,C) = \tfrac12\|\Phi - ( A,B,C)\|_{F}^{2}
            + \lambda_1\sum_r\bigl(\|a_r\|_1+\|b_r\|_1+\|c_r\|_1\bigr),
\end{equation}
optionally replacing $\|\cdot\|_1$ by group-sparse or total-variation terms
\cite{papalexakis_parcube_2012, kim_sparse_2007}, and are solved by
ALS or AO-ADMM.

% ----------------------------------------------------------------------
\subsection{Numerical comparison on silicon FC3}
\label{sub:fc3_results}
% ----------------------------------------------------------------------

We fit the four currently working ans\"atze of \cref{sub:fc3_matsvd,sub:fc3_tucker,sub:fc3_cp,sub:fc3_indscal}
to the mass-weighted, primitive-row FC3 tensor
$\Phi\in\mathbb{R}^{\ndof\times(3\nsc)\times(3\nsc)}$ produced by finite-displacement
DFT in a $3{\times}3{\times}3$ supercell of diamond silicon. The working dimensions are
$\nprim=2$, $\ndof=3\nprim=6$, and $3\nsc=162$, so the dense primitive-row tensor has
\begin{equation}\label{eq:fc3_full_entries}
  N_{\mathrm{full}}
  = \nprim\cdot 3\cdot (3\nsc)^{2}
  = 6\cdot 162^{2}
  = 157\,464
\end{equation}
double-precision entries ($1.26$~MB of raw storage). Thanks to the short-ranged character of the
FC3 interaction in silicon already noted in \cref{sub:fc3_compression}, only a small fraction of
these entries is non-zero, so the natural ``zero-cost'' compression baseline is the number of
non-zero entries $N_{\mathrm{nnz}}$ of the reference tensor, rather than $N_{\mathrm{full}}$
itself.

All transport numbers are obtained by feeding the reconstructed $\tilde\Phi$ back into the same
SCBA + NEGF pipeline used for the dense reference: a single device slab, a $4{\times}4$
transverse $\qp$-mesh, $141$ positive frequency bins in $[0,18]$~THz, $T=300$~K and
$\Delta T=10$~K, Anderson mixing of depth~$5$, SCBA tolerance $10^{-10}$, and at most $50$
SCBA iterations. The reported figure of merit is the anharmonic contribution
$G_{\mathrm{anh}}$ to the thermal conductance per area (MW\,m$^{-2}$\,K$^{-1}$); the
elastic, ballistic conductance is identical across all runs by construction and is not shown.
Errors on both axes are relative to the dense reference, computed once with the raw FC3.

\paragraph{Fit protocol.} All fits are measured by the relative Frobenius error
\begin{equation}\label{eq:fc3_rel_err}
  \varepsilon(\tilde\Phi)
  = \frac{\|\Phi - \tilde\Phi\|_F}{\|\Phi\|_F}
\end{equation}
on the primitive-row tensor. The mSVD and HOSVD ans\"atze are computed in closed form from
LAPACK SVDs of the relevant matricisations; HOSVD is followed by $8$ steps of HOOI refinement.
The CP and INDSCAL fits combine a CP-ALS initialisation (from the \texttt{tensorly} package) with
progressive warm-starts from the previous rank, followed by $400$--$600$ L-BFGS iterations with
a strong-Wolfe line search on the full factor set; $3$--$4$ random restarts are taken per rank
and the best-of-restarts fit is kept. Each ansatz is fitted in two variants: an unconstrained
``plain'' variant and an ``ASR'' variant in which the factor matrices are iteratively projected
onto the acoustic-sum-rule null space during optimisation.

\paragraph{Waring: current status.} The Waring / symmetric-CP fitter of
\cref{sub:fc3_symcp} does not yet converge reliably on this reference tensor and is therefore
omitted from the numerical results below. The present implementation combines shifted
symmetric higher-order power-method (SS-HOPM) initialisation \cite{kolda_shifted_2011} with an
L-BFGS refinement of a single shared factor matrix. On ranks $R\ge 16$ we observe the
classical symptoms of the non-closed rank-$R$ set \cite{stegeman_degeneracy_2006,
de_silva_tensor_2008}: the loss decreases slowly while the factor norms $\|v_r\|$ grow
unboundedly and near-cancelling rank-one terms appear, so that stopping at any finite
iteration count gives a residual that depends sensitively on the random seed and is often
worse than the INDSCAL fit at the same rank. Norm penalties, trust-region L-BFGS, and
degeneracy-aware line searches \cite{stegeman_degeneracy_2006, kolda_numerical_2015} are all
under consideration but are not yet stable enough to include in the sweep.

\paragraph{Frobenius convergence.} \Cref{fig:fc3_frob_and_ganh}~(a) plots $\varepsilon$ versus
parameter count for the four working ans\"atze. At each rank mSVD dominates the
Eckart--Young--Mirsky bound in its chosen bipartition; at the full mode-$3$ rank $R=160$ its
residual collapses to machine precision and the transport observable equals the dense reference
to all reported digits. In the informative regime $\varepsilon\gtrsim 10^{-2}$ we observe
that, \emph{at matched rank}, the accuracy ordering follows the Eckart--Young hierarchy anticipated
in \cref{sub:fc3_compression}: at $R=8$ the Frobenius residuals are
$\varepsilon^{\mathrm{mSVD}}=0.128$, $\varepsilon^{\mathrm{HOSVD}(3,8)}=0.611$,
$\varepsilon^{\mathrm{CP}}=0.164$ and $\varepsilon^{\mathrm{INDSCAL}}=0.166$, with CP and
INDSCAL tracking each other to within $1\%$ relative; at $R=32$ the residuals settle at
$\{0.029,0.033^{(6,32)},0.050,0.052\}$, a factor $\sim 1.7$ gap between mSVD and
CP/INDSCAL and a factor $\sim 10\%$ gap between HOSVD and mSVD once the HOSVD mode-$1$ rank is
saturated at $R_1=6$.

\emph{At matched parameter count} (\cref{fig:fc3_frob_and_ganh}~(a)) the ordering reverses.
With $\nprim=2$, $3\nsc=162$ the leading $R$-scalings read
$N^{\mathrm{mSVD}}\!\approx 1{,}135\,R$, $N^{\mathrm{CP}}\!\approx 331\,R$ and
$N^{\mathrm{INDSCAL}}\!\approx 168\,R$, so a single mSVD term costs $\sim 7\times$ more storage
than a single INDSCAL term. Representative Pareto-frontier points are
\begin{itemize}\setlength{\itemsep}{1pt}
  \item INDSCAL $R=8$: $1\,344$ parameters, $\varepsilon=0.166$ ($117\times$ compression vs.\
        $N_{\mathrm{full}}$);
  \item INDSCAL $R=16$: $2\,688$ parameters, $\varepsilon=0.111$ ($59\times$);
  \item INDSCAL $R=32$: $5\,376$ parameters, $\varepsilon=0.052$ ($29\times$);
  \item CP $R=32$: $10\,592$ parameters, $\varepsilon=0.050$ ($15\times$);
  \item HOSVD $(6,32)$: $8\,388$ parameters, $\varepsilon=0.033$ ($19\times$);
  \item mSVD $R=32$: $36\,320$ parameters, $\varepsilon=0.029$ ($4.3\times$);
  \item INDSCAL $R=256$: $43\,008$ parameters, $\varepsilon=0.012$ ($3.7\times$);
  \item mSVD $R=128$: $145\,280$ parameters, $\varepsilon=8.5\times 10^{-3}$ ($1.1\times$).
\end{itemize}
INDSCAL dominates the Pareto frontier at all parameter counts $\lesssim 3\times 10^{4}$,
followed by HOSVD, CP and finally mSVD. The gap between INDSCAL and CP is modest (a factor
$\sim 2$ in parameters at matched $\varepsilon$) and is entirely driven by the merging of the
two internal factors $b_r\equiv c_r$, which halves the per-rank cost without degrading the
per-rank accuracy on an $S_3$-symmetric tensor; the CP and INDSCAL curves overlap almost
perfectly when plotted against rank.

Two more observations are worth flagging. First, at large ranks the iterative CP and INDSCAL
fits show a clear plateau: the Frobenius residual levels off around $\varepsilon\sim 10^{-2}$
even as $R$ grows from $128$ to $256$, going from $1.5\%$ to $1.2\%$ (INDSCAL) while the
parameter count doubles. This is the classical CP-ALS ``swamp'' \cite{kolda_tensor_2009,
kolda_numerical_2015}, compounded here by the L-BFGS refinement stopping at a local minimum;
mSVD and HOSVD, which have closed-form solutions, do not exhibit it. Second, the ASR-projected
and unconstrained variants produce essentially identical Frobenius curves (differences in the
fourth decimal place, except at $R\le 4$ where the fits are dominated by random initialisation
noise): on an exactly symmetric reference tensor the unconstrained optimum already lies very
close to the ASR-invariant subspace, and the projection costs very little in reconstruction
error. The situation for transport is not symmetric --- see below.

\begin{figure}[t]
  \centering
  \includegraphics[width=\linewidth]{figures/fc3_prim_vasp/frob_and_ganh_vs_params.pdf}
  \caption{FC3 compression on the $3{\times}3{\times}3$ silicon primitive-row tensor
    ($\ndof=6$, $3\nsc=162$, $N_{\mathrm{full}}=157\,464$).
    (a)~Relative Frobenius error \eqref{eq:fc3_rel_err} versus parameter count on log--log
    axes. (b)~Absolute anharmonic conductance $G_{\mathrm{anh}}$ (MW\,m$^{-2}$\,K$^{-1}$)
    versus parameter count on a semi-log $x$ axis. The dashed horizontal line marks the dense
    reference $G_{\mathrm{anh}}^{\mathrm{dense}}=757.46$~MW\,m$^{-2}$\,K$^{-1}$ and the
    shaded band a $\pm 5\%$ corridor; the dotted vertical line marks $N_{\mathrm{nnz}}$, the
    sparse ``zero-cost'' baseline. Solid markers: unconstrained fits. Open markers: ASR-projected
    fits. The mSVD $R=160$ point sits on top of the dense line to machine precision.
    Waring is omitted because the current solver does not converge reliably (see text).}
  \label{fig:fc3_frob_and_ganh}
\end{figure}

\paragraph{Transport convergence and Frobenius--transport decoupling.} The transport axis in
\cref{fig:fc3_frob_and_ganh}~(b) and the left panel of \cref{fig:fc3_kappa_and_spectral}
converges dramatically faster than the Frobenius residual. Already at rank $R=8$ the relative
conductance error $\varepsilon_G = |G - G_{\mathrm{dense}}|/G_{\mathrm{dense}}$ is
$0.78\%$ (mSVD, $\varepsilon=12.8\%$), $0.71\%$ (CP, $\varepsilon=16.4\%$) and $0.33\%$
(INDSCAL, $\varepsilon=16.6\%$); by $R=16$--$32$ all methods sit below $0.5\%$ in transport
while their Frobenius residuals are still in the range $5$--$11\%$. The reduction from $R=4$
to $R=8$ is especially striking: the INDSCAL Frobenius residual improves by a factor $\sim 4$
($0.63 \to 0.17$) while $\varepsilon_G$ improves by a factor $\sim 50$ ($16.9\% \to 0.33\%$).

This decoupling is not a quadratic ``$\varepsilon^{2}$'' effect (which would predict
$\varepsilon_G\sim\varepsilon^{2}\sim 3\%$ at $\varepsilon=0.17$, i.e.\ an order of magnitude
worse than observed). Rather it reflects the strong alignment between the leading tensor
components extracted by the fitters and the phonon-conserving scattering channels that dominate
the conductance integral
\begin{equation}
  G_{\mathrm{anh}} = \int_{0}^{\infty}\!\frac{d\omega}{2\pi}\,\hbar\omega\,J(\omega),
\end{equation}
where $J(\omega)$ is the spectral heat current (\cref{fig:fc3_kappa_and_spectral}~(b)).
Within the compressed ansatz the first few rank-one terms capture the largest singular
components of $\Phi$, which on silicon happen to be aligned with the on-shell three-phonon
processes that satisfy $\omega_a\simeq\omega_c+\omega_d$; the off-shell residual that the
compression discards contributes at high order in $\omega$ to the bubble but makes only a
small contribution to the resonant part of the self-energy that builds up $G_{\mathrm{anh}}$.
In particular, the spectral heat current curves in
\cref{fig:fc3_kappa_and_spectral}~(b) are visually indistinguishable from the dense reference
already at INDSCAL $R=16$--$24$ or mSVD $R=16$, while the corresponding Frobenius residuals are
$10\%$ or more; only at $R\le 4$ does the spectral current show visible distortion, tracking the
Frobenius residual qualitatively.

The practical consequence is that on this system the relevant quality metric for phonon NEGF is
the transport residual $\varepsilon_G$ (or the spectral current residual), not the Frobenius
norm that the fitters themselves minimise. Targeting $\varepsilon_G\le 1\%$ is already achieved
by INDSCAL at $R=8$ ($1\,344$ parameters, $117\times$ compression) and mSVD at $R=8$ ($9\,080$
parameters, $17\times$ compression); reaching $\varepsilon_G\le 0.1\%$ costs INDSCAL $R=32$
($29\times$), HOSVD $(6,32)$ ($19\times$) or mSVD $R=32$ ($4.3\times$). With the current
Waring solver unavailable, INDSCAL is the Pareto-optimal choice across the entire parameter
range we have tested, by one to two orders of magnitude over mSVD.

\begin{figure}[t]
  \centering
  \includegraphics[width=\linewidth]{figures/fc3_prim_vasp/kappa_err_and_spectral.pdf}
  \caption{(a)~Relative thermal-conductance error
    $|G-G_{\mathrm{dense}}|/G_{\mathrm{dense}}$ versus parameter count on log--log axes;
    dotted vertical line: $N_{\mathrm{nnz}}$ reference. Already at $\sim 10^{3}$
    parameters all methods sit below $10^{-2}$, orders of magnitude below their Frobenius
    residual at the same parameter count (\cref{fig:fc3_frob_and_ganh}(a)).
    (b)~Spectral heat current $J(\omega)$ for representative fits against the dense reference
    (red) and the purely ballistic current (dotted grey); the compressed curves overlay the
    dense reference visibly at INDSCAL $R\ge 16$ and mSVD $R\ge 16$.}
  \label{fig:fc3_kappa_and_spectral}
\end{figure}

\paragraph{Conservation of heat current.} The heat-flow conservation residual
$|J_L-J_R|/(|J_L|+|J_R|)$ remains below $0.2\%$ across the entire sweep. The dense reference
sits at $0.01\%$ (limited by the SCBA discretisation floor analysed in
\cref{sub:scba_conservation}); all compressed fits stay within a factor $\sim 20$ of that floor,
with the largest excursions --- still only $1.66\%$ for CP $R=4$ and $1.29\%$ for INDSCAL $R=2$
--- at the very lowest ranks where $\varepsilon$ is $60$--$80\%$. Once $R\ge 8$ the conservation
residual falls back to $0.01$--$0.15\%$ for every method, and shows no systematic trend with
rank. The ASR-projected variants do not improve this floor in any measurable way, consistent
with the observation above that the unconstrained optima are already ASR-compliant to within
fit noise.

\paragraph{Comparison with real-space truncation.} The contrast with a real-space cut-off
baseline is threefold. \emph{Accuracy axis:} on bulk Si a $3$NN pair cut-off fixes
the BTE thermal conductivity within $\sim 1$--$2\%$ of the converged value
\cite{guo_quantum_2020, mcgaughey_phonon_2019}; the same tolerance is reached here by INDSCAL
$R\le 8$ on the NEGF conductance, i.e.\ $1{,}344$ parameters on a reference tensor of
$157{,}464$ entries, and the error can be driven down monotonically to any level by increasing
$R$. A geometric cut-off has no such knob once the supercell is fixed. \emph{Storage axis:}
INDSCAL $R=16$ stores the same tensor at $59\times$ compression with sub-$0.5\%$ transport
error; a $3$NN cut-off on the $3{\times}3{\times}3$ supercell retains the near totality of the
$\sim 1{,}800$ non-zero entries of this particular tensor, i.e.\ gives no additional compression
over the sparse baseline $N_{\mathrm{nnz}}$ (dotted line in
\cref{fig:fc3_frob_and_ganh,fig:fc3_kappa_and_spectral}). \emph{Structure axis:} INDSCAL
delivers a representation that is exactly $S_2$-symmetric in $(j,k)$, so the SCBA kernel
reduces to $\cO(R^{2}\ndof)$ per frequency rather than the $\cO(\ndof^{4})$ of a dense or
cut-off FC3. Matching real-space cut-offs therefore requires either very large supercells
(expensive in displacement DFT) or one of the compressions above, and only the latter
simultaneously accelerates the self-energy contraction.

\subsection{Conservation and SCBA convergence}
\label{sub:scba_conservation}

The accuracy of an FC3 compression is only one axis. A compressed vertex must also
(i) preserve the symmetries that make the SCBA sunset diagram a conserving approximation
in the Baym--Kadanoff sense~\cite{baym_conservation_1961,baym_self-consistent_1962}, and
(ii) keep the resulting self-energy inside the Born-series radius of convergence so that
the fixed-point iteration can be executed at all. We examine both requirements in turn on
the silicon reference tensor.

\paragraph{Baym--Kadanoff conservation of the SCBA sunset.} With the fluctuation vertex
entering the two-loop $\Phi$-functional
\begin{equation}
\Phi[G] = \frac{1}{12}\sum_{abcdef}\Phi_{acd}\,\Phi_{bef}\!
\iint\frac{d\omega_1\,d\omega_2}{(2\pi)^2}\,G_{ab}(\omega_1{+}\omega_2)\,
G_{cf}(\omega_1)\,G_{de}(\omega_2),
\label{eq:lw_functional}
\end{equation}
the SCBA self-energy $\Sigma_S = \delta\Phi/\delta G$ satisfies the integrated continuity
identity $J_S = \tfrac{\hbar}{2}\!\int\!\tfrac{d\omega}{2\pi}\,\omega\,\Tr[\Sigma_S^<G^> -
\Sigma_S^>G^<]\!=\!0$ at the self-consistent fixed point, so that $J_L + J_R = 0$.
Two algebraic properties are required for $\Sigma_S\!=\!\delta\Phi/\delta G$ to close:
full $S_3$ permutation symmetry $\Phi_{acd}\!=\!\Phi_{(acd)}$ of the vertex, which is the
condition under which \cref{eq:lw_functional} is well defined; and the acoustic sum rule
$\sum_{l'\kappa'}\Phi_{0\kappa,l'\kappa',l''\kappa''}=0$, which is the condition under
which the $q\!=\!0$ limit of the bubble vanishes and the $\omega$-integral converges. A
compression scheme that destroys either property loses current conservation, not as a
discretisation artefact but at the algebraic level of the Baym construction.

\paragraph{Numerical verification: perturb $S_3$ and ASR independently.} Starting from the
exactly symmetric silicon FC3 (ASR residual $3.8\times 10^{-15}$ relative to
$\|\Phi\|_F$), we add controlled orthogonal perturbations $E_{S_3}$ ($S_3$-antisymmetric
on the lift, ASR-projected) and $E_{\mathrm{ASR}}$ ($S_3$-symmetric by outer product,
ASR-violating) at magnitudes $\varepsilon\!\in\![10^{-4},0.3]$ and run the full SCBA. The
measured conservation error $|J_L\!-\!J_R|/(|J_L|\!+\!|J_R|)$ sits at the baseline
$6.6\times 10^{-7}$ for $\varepsilon\!\le\!10^{-3}$ on both perturbations (floor-limited,
see below) and separates above $\varepsilon\!\ge\!10^{-2}$, where the ASR break produces a
larger response than the $S_3$ break at matched Frobenius norm. At $\varepsilon\!\sim\!
10^{-1}$ the SCBA enters the non-convergent regime and the scaling becomes non-monotonic.
The correct takeaway is that both symmetries are needed in principle, that the numerical
signature of their violation sits above the SCBA floor, and that the two violations
cannot be separated by perturbative Frobenius magnitude alone on a device of this size.

\paragraph{Decomposition of the discretisation floor.} The $6.6\times 10^{-7}$ baseline
for the exact FC3 is a numerical floor set by the discretised SCBA. Three sources
contribute, which we isolate by varying one knob at a time on a $2{\times}2$ transverse
mesh with $400$ positive frequency bins and Anderson-accelerated SCBA:
(A)~the self-consistency residual $\|\delta\Sigma_S\|$, set by the SCBA tolerance;
(B)~circular-convolution wrap-around of the frequency bubble when $f_{\max}$ is smaller
than twice the phonon bandedge; and (C)~the trapezoidal discretisation of
$\int\!d\omega\,\omega\,\Tr[\Sigma_S^< G^> - \Sigma_S^>G^<]$. On silicon $\omega_{\max}
\!\approx\!16\,\mathrm{THz}$, so $f_{\max}\!\ge\!32\,\mathrm{THz}$ suffices to eliminate
(B) entirely (the edge ratio $|\Sigma^R(f_{\max})|/\max|\Sigma^R|$ drops by four decades
between $f_{\max}\!=\!16$ and $f_{\max}\!=\!48\,\mathrm{THz}$ without moving
$|J_L\!-\!J_R|/(|J_L|\!+\!|J_R|)$). Tightening $\scbatol$ below $10^{-6}$ saturates (A).
What remains -- the dominant contribution to the $6.6\times 10^{-7}$ floor -- is error~(C):
halving $\Delta\omega$ at fixed $f_{\max}\!=\!32\,\mathrm{THz}$ halves the conservation
error, from $9.7\times 10^{-7}$ at $N_\omega\!=\!200$ to $2.0\times 10^{-7}$ at $N_\omega
\!=\!1600$. The $\propto\!\Delta\omega$ scaling (rather than $\Delta\omega^{2}$ of a
smooth integrand) reflects the non-smooth spectral structure of the self-energy at the
phonon band edges, consistent with the Baym proof requiring the integral equalities only
in the continuum limit.

\paragraph{Born-series convergence of the SCBA iteration.} A conserving ansatz is useless
if the fixed-point iteration does not converge. Treating the SCBA update as a map
$(G,\Sigma_S)\!\mapsto\!(G',\Sigma_S')$ through \cref{eq:lw_functional}, Banach contraction
requires the spectral radius $\rho$ of its linearisation at the fixed point to be below
unity. In the phonon case this is the statement of the Neumann-series bound
$\|\Sigma G\|_{\mathrm{op}}\!<\!1$~\cite{vanLeeuwen_Stefanucci_NonequilibriumManyBody_2013},
equivalent at weak coupling to Allen's dimensionless anharmonic parameter
$\alpha_3 = |\Phi_3|^{2}/(\omega^{3}M^{3/2})\!\lesssim\!1$~\cite{allen_anharmonic_2015}
and, at strong coupling, to the Ioffe--Regel bound $\Gamma(\omega)\!\lesssim\!\omega/\pi$
under which the quasiparticle picture survives~\cite{simoncelli_unified_2019,
ravichandran_unified_2020}. We estimate $\rho$ directly by running SCBA with unit linear
mixing ($\alpha\!=\!1$, no Anderson): the asymptotic residual ratio
$\rho = r_{k+1}/r_k$ is the numerical signature of the linearised map's dominant
eigenvalue.

\paragraph{Three criteria isolated on synthetic FC3.} We build random tensors $T$ of
shape $(n_{\mathrm{dof}},\dim_{\mathrm{sc}},\dim_{\mathrm{sc}})$ with a controlled
distance cutoff $r_{\mathrm{cut}}\!\in\!\{2.5,4,6,12\}\,\mathrm{\AA}$ applied on the
$S_3$-symmetric supercell lift (masking pairs $(j,k)$ with
$\max(r_{0j},r_{0k},r_{jk})\!>\!r_{\mathrm{cut}}$) followed by iterative ASR projection,
then rescale to $\alpha \|\Phi_{\mathrm{Si}}\|_F$ with
$\alpha\!\in\!\{0.03,0.1,0.3,1.0\}$. \Cref{fig:scba_convergence} shows the measured
$\rho$ (left) and the Anderson-SCBA conservation error (right) on the resulting
$4{\times}4$ grid. Three regularities emerge.
\emph{First}, ASR is indispensable: a tensor with full $S_3$ symmetry but \emph{without}
ASR projection (not in the figure) produces residuals that oscillate at $\mathcal{O}(1)$
from the first iteration and never enters a basin of attraction, regardless of $\alpha$.
Enforcing ASR drops $\rho$ below $1$ for $\alpha\!\le\!0.1$ across the full range of
$r_{\mathrm{cut}}$. \emph{Second}, coupling is the leading-order knob: $\rho$ grows
essentially exponentially with $\alpha$, crossing unity near $\alpha\!\approx\!0.1$--$0.3$
for every random tensor tested, in line with Allen's criterion applied to
$\|\Phi\|/\omega_{\max}^{2}$. \emph{Third}, the interaction of short-range support with
ASR is non-trivial: narrowing $r_{\mathrm{cut}}$ below a supercell dimension makes the
distance mask and the ASR projector fight each other (they do not commute on a finite
supercell), producing a residual ASR leakage that amplifies $\rho$ at fixed $\alpha$; in
the heatmap this shows as $\rho\!\gtrsim\!1$ for $r_{\mathrm{cut}}\!=\!2.5\,\mathrm{\AA}$
even at $\alpha\!=\!0.1$, while $r_{\mathrm{cut}}\!=\!12\,\mathrm{\AA}$ (no mask) gives
$\rho\!\approx\!10^{-2}$ at the same $\alpha$.

\begin{figure}[t]
  \centering
  \includegraphics[width=0.49\linewidth]{figures/probe_fc3_compactness.pdf}
  \includegraphics[width=0.49\linewidth]{figures/probe_scba_convergence.pdf}
  \caption{Left: cumulative squared Frobenius mass of the FC3 versus maximum
  inter-atomic separation for silicon and for synthetic $S_3$+ASR-symmetric
  tensors of the same shape and norm. Silicon concentrates $\sim\!30\%$ of its mass
  on-site and $\sim\!75\%$ inside $2.5\,\mathrm{\AA}$; random tensors of the same norm
  distribute their mass uniformly out to the supercell boundary. Right: SCBA residual
  histories under pure linear mixing $(\alpha\!=\!1$, no Anderson). Silicon reaches
  $10^{-14}$ before drifting upward past iteration $30$ (marginal fixed point,
  Anderson-stabilisable). Synthetic full-norm $S_3$+ASR tensors converge geometrically
  to $10^{-10}$ then diverge; the same tensor at $0.1\|\Phi_{\mathrm{Si}}\|_F$ contracts
  at ratio $\rho\!\approx\!10^{-2}$; the $S_3$-only tensor (ASR removed) oscillates at
  $\mathcal{O}(1)$ from the outset.}
  \label{fig:scba_convergence}
\end{figure}

\begin{figure}[t]
  \centering
  \includegraphics[width=0.95\linewidth]{figures/probe_range_coupling.pdf}
  \caption{SCBA convergence (left) and Anderson conservation floor (right) for random
  $S_3$+ASR tensors as a function of the distance cutoff $r_{\mathrm{cut}}$ and the
  coupling $\alpha\!=\!\|\Phi\|_F/\|\Phi_{\mathrm{Si}}\|_F$. Contraction ratio
  $\log_{10}\rho$ from $\alpha\!=\!1$ linear-mix SCBA. The white line
  $\rho\!\sim\!1$ crosses the grid diagonally; above it the fixed point is not
  contractive under linear mixing. The real silicon vertex has
  $\alpha\!\equiv\!1$ and an effective $r_{\mathrm{cut}}\!\lesssim\!6\,\mathrm{\AA}$
  yet achieves a conservation floor of $6.6\times 10^{-7}$, three to four decades below
  any random tensor with the same $(r_{\mathrm{cut}},\alpha)$. The gap is the numerical
  signature of the phonon selection rules (energy and crystal-momentum conservation)
  that suppress the three-phonon bubble in the physical vertex but are absent from the
  random construction.}
  \label{fig:range_coupling}
\end{figure}

\paragraph{Silicon as a marginal case.} On the same figure silicon sits at
$\alpha\!=\!1$ and displays a linear-mix $\rho\!\approx\!1.3$: pure linear iteration
diverges, while Anderson acceleration recovers convergence to $10^{-10}$ in $\le\!30$
iterations and a conservation floor of $6.6\times 10^{-7}$, three to four decades below
any synthetic tensor at the same $\alpha$ we could construct. The structural property
responsible is not captured by the Frobenius norm or by the multilinear rank (real and
synthetic have identical mode ranks $[6,45,45]$ at $10^{-4}$ threshold, with ASR
reducing each internal-leg dimension from $48$ to $45$), but by the alignment of
$\Phi_{acd}$ with the harmonic phonon spectrum -- the energy- and
crystal-momentum-conservation selection rules that put the bulk of $|\Phi_{acd}|^{2}$
mass on modes with $\omega_a\!\simeq\!\omega_c\!+\!\omega_d$. This is precisely the
mechanism behind the $3$NN sufficiency reported in real-space cut-off studies of
thermal conductivity~\cite{broido_intrinsic_2007,li_shengbte_2014}, and it is one more
reason why a compression that preserves the algebraic symmetries of $\Phi$ is not the
same thing as one that preserves its \emph{spectral} short-rangedness.

\paragraph{Implications for the compression choice.} Viewed through this lens, the
compression ans\"atze of \cref{sub:fc3_matsvd,sub:fc3_tucker,sub:fc3_cp,sub:fc3_indscal,sub:fc3_symcp}
divide into two classes. mSVD, HOSVD and CP do not enforce $S_3$ symmetry on the vertex
and, unless the residual is driven well below the magnitude of the symmetry-violating
component, they feed the SCBA an effectively anti-symmetric $\Sigma$ that breaks
Baym--Kadanoff closure at the algebraic level. INDSCAL (by construction $S_2$-symmetric)
and Waring (by construction $S_3$-symmetric) preserve the algebraic closure at any rank,
and do not introduce spurious ASR leakage beyond the underlying Frobenius error, because
the translation-invariant lift used during the fit respects the acoustic constraint
through the symmetric factor parametrisation. On the SCBA-convergence axis all five
ans\"atze inherit the Frobenius norm of the input, and none of them tampers with the
coupling $\alpha$ below its reference value. In practice the relevant choice is therefore
between the Pareto-optimal Waring at matched $\varepsilon$ (best storage, exact $S_3$,
degraded fit robustness) and INDSCAL (next-best storage, exact $S_2$, more robust fit);
both remain within the same SCBA regime as the uncompressed silicon reference, with a
conservation floor no worse than the $\varepsilon^{2}$ bound predicted by first-order
Baym--Kadanoff analysis and, at the rank budgets of interest, empirically closer to the
$6.6\times 10^{-7}$ discretisation floor than to any symmetry-violation prediction.